\section{The FeatureLiftBench Benchmark}
\label{sec:benchmark}

\subsection{Repository-Level Feature-Lifting Setting}

A task consists of a pinned source repository $R$, a public contract $S$, an
initial workspace $W$, and a private evaluator $J$. Given $(R,S,W)$, an agent
produces a package $P$ under a fixed execution protocol:
\begin{equation}
  P = \operatorname{Run}(M, H, R, S, W),
  \label{eq:run}
\end{equation}
where $M$ is the model backend and $H$ is the fixed \mainprotocol runtime. The
evaluator observes the collected package and assigns gate-level outcomes:
\begin{equation}
  \operatorname{Functional}(P)
  = B(P) \land P_{\mathrm{pub}}(P)
    \land H_{\mathrm{hid}}(P) \land I(P).
  \label{eq:functional}
\end{equation}
Here, $B$ checks installation and imports, $P_{\mathrm{pub}}$ checks primary
behavior from the public contract, $H_{\mathrm{hid}}$ checks deeper
combinations and boundary behavior from that same contract, and $I$ checks
independence from the source repository. A valid Hidden test must map to a
published contract requirement.

The agent receives the complete source tree and contract but no source-location
hint. It does not see Public tests, Hidden tests, evaluator code, or the frozen
reference. Extraction, adaptation, and behaviorally equivalent reimplementation
are allowed. Runtime imports from the source project, forbidden paths, and
undeclared dependency channels are disallowed.

\begin{figure}[t]
  \centering
  \fbox{%
    \parbox{0.92\linewidth}{\centering
      Pinned repository + public contract\\[2pt]
      $\downarrow$\\[2pt]
      Full-Repository / No-Hint agent execution\\[2pt]
      $\downarrow$\\[2pt]
      Standalone \texttt{featurelifted} package + trajectory\\[2pt]
      $\downarrow$\\[2pt]
      Build $\rightarrow$ Public $\rightarrow$ Hidden $\rightarrow$ Isolation\\[2pt]
      $\downarrow$\\[2pt]
      Functional Pass + pass-conditioned compactness}}
  \caption{The \flb evaluation pipeline. Protected tests, evaluator code, and
  evaluator-side references are never exposed during agent execution.}
  \Description{A pipeline from a pinned repository and public contract through
  agent execution and a standalone package to build, Public, Hidden, and
  isolation evaluation, followed by functional and compactness reporting.}
  \label{fig:pipeline}
\end{figure}

\subsection{Task Suite Design}

The paper suite, Python-200$'$, contains 200 tasks from 176 repositories
(Table~\ref{tab:suite}). It emphasizes Python libraries, developer tools, and
framework components. One source repository may support multiple distinct
tasks in the baseline split. Every task pins its upstream revision and records
source identity in a canonical registry.

\begin{table}[t]
  \caption{Composition of the Python-200$'$ paper suite.}
  \label{tab:suite}
  \centering
  \begin{tabular}{lrrl}
    \toprule
    \textbf{Split} & \textbf{Tasks} & \textbf{Repositories} & \textbf{Role} \\
    \midrule
    Frozen Python-150 & 150 & 127 & Baseline set \\
    Hard-50 & 50 & 50 & Calibrated expansion \\
    \midrule
    Python-200$'$ & 200 & 176 & Main paper suite \\
    \bottomrule
  \end{tabular}
\end{table}

Tasks are organized along three descriptive axes. \emph{Feature family}
records requested behavior, such as parsing, serialization, configuration,
validation, resource handling, registry behavior, caching, protocols,
algorithms, or workflows. \emph{Primary entanglement} records the dominant
mechanism that makes independent extraction difficult: data-model,
parser-state, framework, configuration/environment, resource, or third-party
dependency entanglement. \emph{Lift type} describes the relationship between
the requested artifact and source implementation: Direct, Adapted, or
Composite. These labels support sampling and analysis; they are not shown to
the agent and do not affect Functional Pass.

The combined suite contains 68 Direct, 100 Adapted, and 32 Composite tasks.
The Hard-50 expansion increases coverage of registry/plugin dispatch, workflow
and session orchestration, configuration discovery, validation boundaries,
deep parsing, and direct tooling copy traps. It was calibrated with a strong
model at 29/50 functional passes, within a predeclared 40--65\% target band.
This is benchmark-design evidence, not a substitute for the final main-table
evaluation.

An earlier External-50 expansion is retained as a side split rather than part
of Python-200$'$. Strong-model pass rates of 90--94\% and pass-conditioned
footprints close to one reference-relative unit revealed an easy, copy-heavy
construction. This negative design result motivated Hard-50 and illustrates why
functional outcome and extraction compactness must be measured separately.

\subsection{Task Construction and Validation}

Each candidate task is retained only when it satisfies four requirements:
\begin{enumerate}
  \item \textbf{Realism:} the requested feature corresponds to a plausible
        reuse or modularization need in a real repository;
  \item \textbf{Solvability:} a feasible independent implementation can be
        produced from the pinned repository and public contract;
  \item \textbf{Oracle-checkability:} success can be verified with
        deterministic build, behavior, and isolation checks;
  \item \textbf{Integrity:} the agent cannot obtain credit by reading
        protected tests, retaining a forbidden source dependency, or bypassing
        the intended artifact boundary.
\end{enumerate}

Task construction records a source revision and digest, public contract,
required API, locked dependencies, Public/Hidden contract mapping,
deterministic evaluator, and a feasible reference. Evaluator-side references
are never exposed to agents. Automated and manual gates check source
materialization, package imports, reference behavior, test isolation,
forbidden dependencies, resource access, and freeze identity.

The benchmark separates task packages from mutable working-tree state. Paper
results must be evaluated against the frozen materialization identified by the
active suite manifest. A source mismatch, unavailable locked dependency,
evaluator defect, or context-policy violation is an eligibility or
infrastructure outcome rather than a model failure.

\subsection{Run Protocol and Evidence Collection}

Each run follows a setup--execution--evaluation pipeline
(Figure~\ref{fig:pipeline}). Setup materializes the pinned task workspace and
records task, source, model, runtime, prompt, budget, and container identity.
During execution, the agent may inspect the full repository, use shell and
editing tools, run legal self-tests, and modify its submission workspace. The
runner collects the final submission tree, model usage, actions, errors, and
workspace snapshots. Evaluation installs the submission in a Dockerized
offline environment and applies build, Public, Hidden, and isolation checks.

Each task attempt produces four evidence layers: (1) the final artifact,
(2) the agent trajectory and process status, (3) token, step, time, context,
and resource statistics, and (4) the task-level deterministic evaluator result.
The evaluator result is authoritative for Functional Pass. Agent completion or
runner status is reported separately because an agent may hit a step limit after
leaving a passing package, or may report completion while its artifact fails.

\subsection{Metrics}

\paragraph{Functional Pass@1.}
A task receives one functional point only if all gates in
Equation~\ref{eq:functional} pass. Functional failures receive one exclusive
first-failure stage:
\begin{equation}
  \text{missing} \rightarrow \text{build} \rightarrow \text{public}
  \rightarrow \text{hidden} \rightarrow \text{isolation}.
\end{equation}
Infrastructure and evidence-availability flags are reported separately from
semantic model outcomes.

\paragraph{Compactness.}
For a functionally passing submission $P$ with evaluator-side feasible
reference $P_{\mathrm{ref}}$, we report
\begin{equation}
  \operatorname{RRES}(P)
  = \frac{\operatorname{normalized\_size}(P)}
         {\operatorname{normalized\_size}(P_{\mathrm{ref}})}.
  \label{eq:rres}
\end{equation}
\rres is accompanied by file, copied-code, and dependency-footprint
diagnostics where evidence is available. It is reported only on functional
passes. Cross-method compactness comparisons use the same-task subset on which
both methods pass.

\paragraph{Process and cost.}
Tokens, steps, latency, completion status, and resource failures are diagnostic
metrics. On passing trajectories with replayable package snapshots, $T^*$ is
the earliest cumulative token point at which a unique package tree passes the
evaluator. $T^*$ is undefined for failing trajectories and is not itself a
deployable stopping rule.
